\section{Evaluation}\label{sec:eval}

We conduct extensive evaluations on \sys and compare it against both open-weight and fully open models. Our \sys advances the frontier of fully open models and further narrows the gap between fully open models and leading open-weight models.

\subsection{Evaluation Setup}

\subsubsection{Baseline Models}
We evaluate \sys against a comprehensive set of state-of-the-art baseline models with comparable parameter counts. These baselines are categorized into two distinct groups: \textbf{open-weight models}, where model weights are public but training data or details may remain proprietary, and \textbf{fully-open models}, where the architecture, weights, training code, and datasets are all publicly released.\footnote{All evaluated models are base (pretrained) checkpoints without instruction tuning. To maintain consistency, we standardize naming conventions by omitting suffixes (e.g., using simplified names for Qwen3) to denote base models throughout this paper.}

\paragraph{Open-weight models.}
\begin{itemize}
    \item \textit{Qwen2-1.5B}~\cite{qwen2}: A 1.5B-parameter decoder-only transformer trained on large-scale multilingual and code data. It delivers robust performance in general language understanding, coding, and reasoning while facilitating efficient deployment.
    \item \textit{Qwen2.5 series}~\cite{qwen2025qwen25technicalreport}: We select Qwen2.5-1.5B and Qwen2.5-3B, dense foundation models that refine the Qwen2.5 architecture. These models feature an improved tokenizer and offer enhanced capabilities in knowledge, coding, and mathematics within a compact form factor suitable for edge applications.
    \item \textit{Qwen3 series}~\cite{qwen3}: We include Qwen3-0.6B, Qwen3-1.7B, and Qwen3-4B. These are small-scale base models that support long contexts and ``thinking'' modes, providing competitive abilities in general tasks, mathematics, and coding.
    \item \textit{Gemma2-2B}~\cite{morgane2024gemma2}: The smallest member of Google's Gemma 2 family, this model is distilled from larger counterparts. It was trained on 2 trillion tokens from diverse sources, including web documents, code, and scientific articles.
    \item \textit{Llama3.2 series}~\cite{llama3.2}: We utilize Llama-3.2-1B and Llama-3.2-3B, multilingual text-only models distilled and pruned from larger Llama variants. They support extended context windows (128k) and tool-calling, targeting privacy-preserving on-device inference.
\end{itemize}

\paragraph{Fully-open models.}
\begin{itemize}
    \item \textit{SmolLM2-1.7B}~\cite{smol2}: Developed by Hugging Face, this model utilizes the Llama 2 architecture with a GPT-2 tokenizer (vocabulary size 49,152). It was trained on 256 H100 GPUs.
    \item \textit{SmolLM3-3B}~\cite{smollm3}: A compact, fully-open model trained on 11T tokens using data-centric recipes. It features a 128k context window utilizing NoPE and YaRN, offering state-of-the-art performance for its size class with multilingual support.
    \item \textit{OLMo2-1B}~\cite{olmo20252olmo2furious}: The smallest model in the OLMo2 family (specifically OLMo2-0425-1B), trained on the OLMo-mix corpus. Its full release of code, checkpoints, logs, and training details enables rigorous scientific inquiry into compute-efficient training at the 1B-parameter scale.
    \item \textit{YuLan-Mini}~\cite{yiwen-etal-2025-yulan}: A 2.4B-parameter model pre-trained on approximately 1.08T tokens. By combining curated data pipelines with robust optimization and annealing strategies, it achieves top-tier performance among similarly sized models, particularly in mathematics and coding.
\end{itemize}

\subsubsection{Benchmarks}
Our evaluation encompasses four primary domains: mathematics, coding, Chinese language processing, and general reasoning \& knowledge. We selected representative benchmarks for each domain as follows:
\begin{itemize}
    \item \textit{Math}: We utilize GSM8K~\cite{gsm8k} and MATH~\cite{MATH}. Together, these datasets cover the spectrum from grade-school arithmetic to advanced competition-style problems, providing a comprehensive assessment of symbolic and multi-step reasoning.
    \item \textit{Coding}: We adopt MBPP~\cite{mbpp} and HumanEval~\cite{humaneval} to evaluate code generation via unit testing. This approach directly measures the model's ability to synthesize executable and logically coherent programs. Specifically, we use the sanitized subset of MBPP, which refines problem descriptions and test cases to minimize ambiguity.
    \item \textit{Chinese}: To assess knowledge and reasoning within a Chinese linguistic context, we employ CMMLU~\cite{cmmlu} and C-Eval~\cite{ceval}, widely adopted benchmarks spanning diverse academic and professional subjects.
    \item \textit{Reasoning \& Knowledge}: For general English-language knowledge and reasoning, we include a suite of eight datasets: MMLU~\cite{mmlu}, HellaSwag~\cite{hellaswag}, Common Sense QA (CSQA)~\cite{commonsenseqa}, BoolQ~\cite{clark-etal-2019-boolq}, PIQA~\cite{piqa}, SocialIQA~\cite{socialiqa}, WinoGrande~\cite{winogrande}, and ARC~\cite{arc}. These benchmarks cover a broad range of expert knowledge, reading comprehension, and commonsense reasoning scenarios.
\end{itemize}

\subsubsection{Implementation Details}
We conduct our evaluation using the OpenCompass framework~\cite{2023opencompass}, a comprehensive platform for large model evaluation. For mathematics and coding benchmarks, which typically require open-ended generation, we evaluate models in \textit{generation mode}. Conversely, for benchmarks in other domains, we employ \textit{perplexity-based (PPL) evaluation}. Following the OLMES protocol~\cite{gu2025olmes}, PPL tasks are assessed under both multiple-choice formulation (MCF) and completion formulation (CF), with the superior score reported as the final result.

\subsection{Evaluation Results}

The performance of \sys and baseline models is summarized in \Cref{tab:lang_math_code,tab:reasoning_knowledge}, with a comprehensive comparison provided in \Cref{tab:model_comparison_full}.

\begin{table}[!htbp]
\centering
\caption{Core Capabilities: Chinese, Math, and Code.}
\label{tab:lang_math_code}
\resizebox{\textwidth}{!}{%
\begin{threeparttable}
\begin{tabular}{lccccccccc}
\toprule
\multirow{3}{*}{\textbf{Model Name}} & \multirow{3}{*}{\textbf{Params}} & \multicolumn{2}{c}{\textbf{Chinese}} & \multicolumn{2}{c}{\textbf{Math}} & \multicolumn{2}{c}{\textbf{Code}} & \multirow{3}{*}{\textbf{Avg}}\\
\cmidrule(lr){3-4} \cmidrule(lr){5-6} \cmidrule(lr){7-8}
 &  & C-Eval & CMMLU  & GSM8K & MATH & sanitized-MBPP & HumanEval & \\
 &  & 5 shot & 5 shot & 4 shot & 4 shot & 3 shot & 3 shot & \\
\midrule

\multicolumn{9}{l}{\textit{\textbf{Open-Weight SOTA}}} \\
\rowcolor{openweight} Qwen2-1.5B & 1.5B & 71.29 & 70.62 & 58.50$^{*}$ & 21.70$^{*}$ & 50.58 & 31.10$^{*}$ & 50.63\\
\rowcolor{openweight} Qwen2.5-1.5B & 1.5B & 68.63 & 68.01  & 68.50$^{*}$ & 35.00$^{*}$ & 58.37 & 37.20$^{*}$ & 55.95\\
\rowcolor{openweight} Qwen2.5-3B & 3B & 74.65 & 73.92  & 79.10$^{*}$ & 42.60$^{*}$ & 66.54 & 42.10$^{*}$ & 63.15\\
\rowcolor{openweight} Qwen3-0.6B & 0.6B & 57.03 & 52.36 & 59.59$^{*}$ & 32.44$^{*}$ & 51.75	 & 29.88 & 47.18 \\
\rowcolor{openweight} Qwen3-1.7B & 1.7B & 66.70 & 66.55  & 75.44$^{*}$ & 43.5$^{*}$ & 64.20 & 52.44 & 61.47\\
\rowcolor{openweight} Qwen3-4B & 4B & 78.5 & 77.01 & 87.79$^{*}$ & 54.10$^{*}$ & 74.32 & 62.20 & 72.32\\
\rowcolor{openweight} Gemma2-2B & 2B & 41.35 & 39.63 & 23.90$^{*}$ & 15.00$^{*}$ & 38.91 & 17.70$^{*}$ & 29.42\\
\rowcolor{openweight} Llama-3.2-1B & 1B & 29.82 & 31.03 & 44.40$^{*}$ & 30.60$^{*}$ & 34.63 & 18.90 & 31.56\\
\rowcolor{openweight} Llama-3.2-3B & 3B & 45.67 & 44.33 & 77.70$^{*}$ & 48.00$^{*}$ & 49.42 & 29.88 & 49.17\\

\multicolumn{9}{l}{\textit{\textbf{Fully-Open SOTA}}} \\
\rowcolor{fullyopen} SmolLM2-1.7B & 1.7B & 35.06 & 34.03 & 31.10$^{*}$ & 11.60$^{*}$ & 49.42 & 22.60$^{*}$ & 30.64\\
\rowcolor{fullyopen} OLMo-2-0425-1B & 1B & 30.53 & 28.62 & 68.30$^{*}$ & 20.70$^{*}$ & 15.56 & 6.71 & 28.40\\
\rowcolor{fullyopen} YuLan-Mini-2.4B & 2.4B & 52.32 & 48.14 & 66.65$^{*}$ & 27.12 & 62.26 & 61.60$^{*}$ & 53.02\\
\rowcolor{fullyopen} SmolLM3-3B & 3B & 50.84 & 49.35 & 67.63$^{*}$ & 46.10$^{*}$	& 62.26 & 39.63 & 52.64\\

\multicolumn{9}{l}{\textit{\textbf{Ours}}} \\
\rowcolor{ours} PCMind-2.1-Kaiyuan-2B & 2B & 46.30 & 49.25 & 51.33 & 30.34 & 56.42 & 42.68 & 46.05\\

\bottomrule
\end{tabular}%

\begin{tablenotes}
\item[*] This score is cited from the corresponding official report or paper.
\end{tablenotes}
\end{threeparttable}
}
\end{table}

\textbf{Core Capabilities: Math, Code, and Chinese.} 
In \Cref{tab:lang_math_code}, we focus on three specialized capabilities of the model: mathematics, coding, and Chinese language\footnote{For generation tasks (math and code), we report official results for baseline models where available, as exact reproduction can be challenging.}. \sys achieves an average score of \texttt{46.05} across these seven benchmarks. It outperforms fully-open models of similar scale, such as SmolLM2-1.7B and OLMo-2-0425-1B, and remains competitive with larger models like YuLan-Mini-2.4B and SmolLM3-3B, despite a smaller parameter count.
Specifically, on Chinese benchmarks (C-Eval and CMMLU), \sys scores \texttt{46.30} and \texttt{49.25}, respectively—markedly higher than SmolLM2-1.7B and OLMo-2-0425-1B, and approaching the performance of the larger SmolLM3-3B. On mathematics, \sys achieves \texttt{51.33} on GSM8K, substantially surpassing SmolLM2-1.7B, and scores \texttt{30.34} on MATH, outperforming YuLan-Mini-2.4B (\texttt{27.12}). Similarly, in code generation, our model reaches \texttt{42.68} on HumanEval, exceeding both SmolLM3-3B (\texttt{39.63}) and Qwen2.5-3B (\texttt{42.10}). The results demonstrate \sys offers a superior trade-off between accuracy and model size in critical domains.

\textbf{Reasoning and Knowledge.}
\Cref{tab:reasoning_knowledge} presents performance on nine reasoning and knowledge benchmarks. \sys achieves an average score of \texttt{67.74}, placing it firmly within the competitive range for its size class. Within the fully-open category, our model surpasses SmolLM2-1.7B (+1.69 average) and OLMo-2-0425-1B (+5.68 average), while effectively matching the larger YuLan-Mini-2.4B (\texttt{67.50}). Although the larger SmolLM3-3B attains a higher average (\texttt{72.60}), \sys significantly narrows the gap to the state-of-the-art for fully-open models at this scale.
When compared to open-weight models, \sys achieves an average score of \texttt{67.74}, approaching Gemma2-2B (\texttt{69.16}) despite using comparable training data. 
Larger open-weight models like Qwen3-4B maintain a substantial lead (\texttt{81.84}), which is expected given their significantly larger scale and training resources.

\begin{table}[htbp]
\centering
\caption{Reasoning and Knowledge Capabilities.}
\label{tab:reasoning_knowledge}
\resizebox{\textwidth}{!}{%
\begin{tabular}{lcccccccccccc}
\toprule
\multirow{3}{*}{\textbf{Model Name}} & \multirow{3}{*}{\textbf{Params}} & \multicolumn{9}{c}{\textbf{Reasoning \& Knowledge}} & \multirow{3}{*}{\textbf{Avg}}\\
\cmidrule(lr){3-11}
 &  & MMLU & ARC-C & ARC-E & BoolQ & CSQA & HSwag & PIQA & SocIQ & Wino & \\
&  & 5 shot & 5 shot & 5 shot & 5 shot & 5 shot & 5 shot & 5 shot & 5 shot & 5 shot &\\
\midrule

\multicolumn{11}{l}{\textit{\textbf{Open-Weight SOTA}}} \\
\rowcolor{openweight} Qwen2-1.5B & 1.5B & 56.36 & 70.17 & 83.60 & 71.90 & 70.52 & 60.77 & 75.73 & 63.46 & 59.83 & 68.04\\
\rowcolor{openweight} Qwen2.5-1.5B & 1.5B & 61.56 & 79.32 & 90.48 & 76.39 & 75.10 & 64.18 & 76.17 & 64.94 & 59.67 & 71.98 \\
\rowcolor{openweight} Qwen2.5-3B & 3B & 66.86 & 86.44 & 92.59 & 83.88 & 76.09 & 73.85 & 81.45 & 69.40 & 63.69 & 77.14\\
\rowcolor{openweight} Qwen3-0.6B & 0.6B & 55.09 & 68.14 & 84.48 & 69.05 & 61.18 & 48.51 & 69.97 & 61.51 & 55.64 & 63.73 \\
\rowcolor{openweight} Qwen3-1.7B & 1.7B & 65.35 & 80.34 & 91.89 & 79.82 & 74.61 & 60.76 & 77.20 & 68.58 & 59.27 & 73.09\\
\rowcolor{openweight} Qwen3-4B & 4B & 75.78 & 89.83 & 97.53 & 86.09 & 81.9 & 79.46 & 84.98 & 75.59 & 65.43 & 81.84\\
\rowcolor{openweight} Gemma2-2B & 2B & 55.20 & 66.44 & 82.54 & 72.42 & 69.45 & 66.20 & 78.89 & 65.92 & 65.35 & 69.16\\
\rowcolor{openweight} Llama-3.2-1B & 1B & 37.74 & 36.95 & 70.55 & 67.43 & 62.82 & 60.20 & 74.92 & 50.61 & 58.17 & 57.71\\
\rowcolor{openweight} Llama-3.2-3B & 3B & 57.87 & 72.20 & 83.95 & 76.73 & 70.35 & 71.06 & 79.05 & 64.33 & 64.09 & 71.07\\

\multicolumn{11}{l}{\textit{\textbf{Fully-Open SOTA}}} \\
\rowcolor{fullyopen} SmolLM2-1.7B & 1.7B & 51.99 & 59.66 & 82.72 & 69.85 & 67.16 & 65.30 & 78.51 & 60.18 & 59.12 & 66.05\\
\rowcolor{fullyopen} OLMo-2-0425-1B & 1B & 44.25 & 47.46 & 76.72 & 70.55 & 65.60 & 61.61 & 76.44 & 55.53 & 60.38 & 62.06\\
\rowcolor{fullyopen} YuLan-Mini-2.4B & 2.4B & 51.76 & 64.75 & 82.54 & 78.59 & 66.18 & 61.20 & 77.31 & 63.25 & 61.88 & 67.50\\
\rowcolor{fullyopen} SmolLM3-3B & 3B & 63.04 & 77.29 & 88.54 & 76.12 & 70.52 & 69.20 & 79.05 & 65.25 & 64.40 & 72.60\\

\multicolumn{11}{l}{\textit{\textbf{Ours}}} \\
\rowcolor{ours} PCMind-2.1-Kaiyuan-2B & 2B & 53.90 & 66.10 & 82.89 & 78.53 & 67.40 & 58.13 & 74.37 & 62.59 & 65.75 & 67.74\\

\bottomrule
\end{tabular}%

}
\end{table}

\paragraph{Discussion on Size and Performance Trade-offs.} The overall trade-off between model size and average benchmark performance is visualized in \Cref{fig:model_perf_comp}. The figure reveals that \sys lies beyond the current fully-open frontier: at comparable parameter counts, it clearly outperforms earlier fully-open models (e.g., OLMo-2-1B, SmolLM2-1.7B) and approaches the performance of the larger YuLan-Mini-2.4B. Moreover, if adhering to the convention of comparing non-embedding parameters to get rid of the vocabulary effect, our \sys can exhibit even more prominent advantages, as shown in \Cref{fig:performance_comparison_on_nonembedding}.
We compare different models according to non-embedding parameters in \Cref{sec:non_embedding_comparison}.

Furthermore, when compared to open-weight baselines of similar size, \sys demonstrates superior architectural efficiency. 
For instance, compared to Gemma2-2B, \sys utilizes fewer model parameters (1.4B non-embedding and 2B total parameters, versus Gemma2-2B's 2B non-embedding and 2.6B total parameters) while maintaining a similar token budget (2.2T tokens for\sys versus 2T for Gemma2-2B). 
Despite these reduced resource requirements, \sys achieves stronger performance on core capabilities (Chinese, Math, Code) and competitive reasoning scores, as shown in \Cref{tab:lang_math_code,tab:reasoning_knowledge,tab:model_comparison_full}. 
Although a performance gap remains compared to the Qwen series, likely attributable to their massive training data scale (e.g., 36T tokens), \sys occupies a favorable position in the size-performance landscape, offering a strong, fully-open alternative for resource-constrained environments.


















